% !TEX program = xelatex
\documentclass[UTF8,a4paper,11pt,fontset=fandol]{ctexart}
\usepackage[margin=26mm]{geometry}
\usepackage{amsmath,amssymb,amsthm,mathtools,booktabs,longtable,array,graphicx,enumitem}
\usepackage{xcolor}
\usepackage[colorlinks=true,linkcolor=blue!50!black,citecolor=blue!50!black,urlcolor=blue!50!black]{hyperref}
\allowdisplaybreaks[2]
\setlength{\emergencystretch}{2em}
\linespread{1.12}
\newtheorem{theorem}{定理}
\newtheorem{oldtheorem}{定理}
\renewcommand{\theoldtheorem}{\Alph{oldtheorem}}
\newtheorem{proposition}[theorem]{命题}
\newtheorem{lemma}[theorem]{引理}
\newtheorem{corollary}[theorem]{推论}
\theoremstyle{definition}\newtheorem{definition}[theorem]{定义}
\theoremstyle{remark}\newtheorem{remark}[theorem]{注}
\newcommand{\R}{\mathbb{R}}
\newcommand{\N}{\mathbb{N}}
\newcommand{\Z}{\mathbb{Z}}
\newcommand{\dd}{\,\mathrm d}
\newcommand{\ind}{\mathbf{1}}
\DeclareMathOperator{\supp}{supp}
\DeclareMathOperator{\essinf}{ess\,inf}
\DeclareMathOperator{\rank}{rank}
\DeclareMathOperator{\spanop}{span}
\DeclareMathOperator{\Rea}{Re}
\title{八维中心方体极大算子的弱型常数\\双侧界与有限方向势方法}
\author{}
\date{2026年10月27日}
\begin{document}
\maketitle
\begin{abstract}
设 $C_8$ 为 Lebesgue 测度下八维中心轴平行方体极大算子的最优弱 $(1,1)$ 常数，所有坐标使用同一半径并取遍全部正尺度。本文证明
\[
 \frac{2188811179426669}{10^{15}}<C_8<
 \frac{56627260639}{2000000000}.
\]
上界来自带权选择法及方体壳上捕获区域的精确估计；下界由非等距格点构造、完整见证盒并集的有理计算及有限光滑输入实现。另对八维有限方向 Cauchy 半群证明 Green 势的奇异集恰为秩不超过三的方向子空间之并，并建立正部势的全空间概率源与卷积比较原理。面上的短时间 $L^{8/7}$ 估计控制这些子空间附近的积分，为进一步构造上界核提供可检验条件。
\par\smallskip\noindent\textbf{关键词：}中心极大算子；弱型常数；轴平行方体；计算辅助证明；Cauchy 半群；Green 势。
\end{abstract}

\tableofcontents
\section{引言}
中心 Hardy--Littlewood 极大算子的弱 $(1,1)$ 不等式将局部平均的几何结构与水平集的体积联系起来。本文固定 Lebesgue 测度及轴平行方体
\[
 Q(x,r)=x+(-r,r)^n,\qquad
 M_nf(x)=\sup_{r>0}\frac{1}{(2r)^n}\int_{Q(x,r)}|f(y)|\dd y.
\]
同一个 $r$ 用于全部坐标。记 $C_n$ 为满足
\begin{equation}\label{eq:weak-definition}
 \lambda\bigl|\{M_nf>\lambda\}\bigr|\le C_n\|f\|_1
 \quad(f\in L^1(\R^n),\ \lambda>0)
\end{equation}
的最小常数。研究固定维数的数值界，需要同时保留共同半径以及一般输入这两个条件：分别估计一维平均的乘积，或只在有限尺度上取极大，都会改变所讨论的问题。

Aldaz 的格点方法给出每个维数的显式下界\cite{Aldaz2000}。令
\[
 u_n=2^{1/n},\quad v_n=(3/2)^{1/n},\quad a_n=(1+u_n)/2,
\]
其 Proposition 1.4 所给的量为
\begin{equation}\label{eq:Aldaz-bound}
 A_n=\frac32a_n^{n-1}
 -\frac12\left[(3a_n-2v_n)^{n-1}
 -\left(\frac{3-u_n}{2}\right)^{n-1}\right].
\end{equation}
特别地 $A_8\approx1.78390524162816$。另一方面，中心 $\ell_\infty$ 球的通用覆盖估计给出 $C_n\le2^n$\cite{AldazKissing2020}，从而 $C_8\le256$。Stein--Str\"omberg 的选择法建立了一般对称凸体的 $O(n\log n)$ 弱型估计\cite{SteinStromberg1983}；在固定维数问题中，还需具体计算覆盖与重叠估计中的常数。Aldaz 的发散定理及后续定量增长结果\cite{Aldaz2011,IakovlevStromberg2013}说明方体常数不能取维数一致的上界，但渐近增长与八维显式数值是两个不同层次的问题。

本文在八维同时构造较大的可实现水平集和较小的普遍覆盖常数。上界的核心是：选择过程留下的未捕获部分位于一个显式方体角帽区域；对坐标壳积分，可从粗重叠估计中扣除始终被捕获的部分。下界的核心是：固定七个间距后，首轴间距对应的完整见证体积连续分段仿射，因此可以按有限端点事件寻找构造，再用准确并集算法认证单个构造。

\begin{theorem}\label{thm:main}
对于上述八维中心轴平行方体极大算子，
\begin{equation}\label{eq:main-bound}
 \boxed{\quad
 \frac{2188811179426669}{10^{15}}<C_8<
 \frac{56627260639}{2000000000}\quad}.
\end{equation}
下界由有限、非负、光滑紧支撑的 $L^1$ 输入实现为一个严格较大的弱型比值；上界对所有 $L^1$ 输入与全部正半径成立。
\end{theorem}

式\eqref{eq:main-bound}的显示形式为
$2.188811179426669<C_8<28.3136303195$。
这里与式\eqref{eq:Aldaz-bound}和通用上界 $256$ 作具体比较，不将这两个比较对象称为固定 Lebesgue 背景下的最优界。计算辅助部分的完整根隔离、并集判据与准确分数作为补充材料给出；主文证明有限判据怎样产生允许的 $L^1$ 输入，并说明它们如何进入定理。

我们还研究一种可用于进一步改进上界的方向势方法。Lin 的正核比较原理\cite{Lin2026}将极大算子的弱型估计归约为具有适当源正性的可积主核。对八维有限方向 Cauchy 过程，Green 势的奇异集已经不只包含方向线：秩不超过三的方向子空间都可能产生不可积的小时间点值。本文证明这一奇异集的准确描述，并以全空间分布方法构造正部势的概率源。由此得到一般输入与全尺度的比较定理，以及控制奇异邻域的面上短时范数估计。这个分析结果与定理\ref{thm:main}的几何上界是两条独立论证。

第二节统一记号和所用标准工具。第三节分别证明几何上界、计算辅助下界与方向势的比较定理。第四节给出强 $L^p$ 推论及十六方向模型的具体计算。附录列出准确证书、验证程序和复现方法。

\section{预备知识}\subsection{测度、尺度与水平集}
全文的体积 $|E|$ 和积分均相对于 Lebesgue 测度。开、闭方体给相同的 $L^1$ 平均；在原子构造中使用闭方体，仅用于描述有限见证，最终输入通过同质量平滑得到。对复值输入，算子作用于 $|f|$。

\begin{lemma}\label{pre:continuity}
若 $f\in L^1(\R^n)$，则每个固定半径的中心方体平均关于中心连续。因此 $\{M_nf>\lambda\}$ 为开集；其测度可由紧子集内逼近。
\end{lemma}
\begin{proof}
平移 $h$ 所引起的平均差不超过
$(2r)^{-n}\||f|(\cdot+h)-|f|\|_1$，它由 $L^1$ 平移连续性趋于零。任意族连续函数的上确界是下半连续函数，结论成立。
\end{proof}

下文采用 Tonelli 定理处理非负积分，采用具有明确支配函数的支配收敛处理有符号极限。Fourier 变换约定为
\[
 \widehat f(\xi)=\int_{\R^n}e^{-ix\cdot\xi}f(x)\dd x,
 \qquad f(x)=(2\pi)^{-n}\int_{\R^n}e^{ix\cdot\xi}\widehat f(\xi)\dd\xi.
\]
涉及 $L^2$ 能量时可等价采用 unitary 归一化。所用 Fourier 反演和 Plancherel 定理取通常的 $L^1$ 及 $L^2$ 版本\cite{Folland1999}；局部 Rellich 紧性取有界球上 $H^1\hookrightarrow L^2$ 紧嵌入版本\cite{Evans2010}。第三节在每次使用处说明相应支配、域条件和紧性来源。

\subsection{已有比较与维数承接}
\begin{oldtheorem}[Aldaz，Proposition 1.4]\label{pre:old-lower}
以式\eqref{eq:Aldaz-bound}定义 $A_n$。对 $n\ge2$，有 $C_n\ge A_n$\cite{Aldaz2000}。
\end{oldtheorem}
\begin{oldtheorem}[中心 $\ell_\infty$ 球的覆盖估计]\label{pre:old-upper}
对每个正整数 $n$，$C_n\le2^n$\cite[Corollary 3.5]{AldazKissing2020}。
\end{oldtheorem}
定理 B 所在文献讨论的测度一致锐性涉及对背景测度取上确界，专门化到固定 Lebesgue 测度只给这里的上界。

\begin{lemma}\label{pre:dimension}
对中心轴平行、共同半径的方体，有 $C_{n+1}\ge C_n$。
\end{lemma}
\begin{proof}
取 $b<C_n$，选择 $f\in L^1(\R^n)$、$\lambda>0$ 及紧集
$K\subset\{M_nf>\lambda\}$ 使 $\lambda|K|>b\|f\|_1$。
引理\ref{pre:continuity}及有限子覆盖给出有限个实际见证半径，令它们的最大值为 $R$。
对 $T>R$，令 $F_T(x,s)=f(x)\ind_{[-T,T]}(s)$。
当 $x\in K$、$|s|\le T-R$ 时，原见证半径 $r\le R$ 在新增坐标的积分因子恰为 $2r$，故相应 $(n+1)$ 维平均等于原 $n$ 维平均。
因此
\[
 \frac{\lambda|\{M_{n+1}F_T>\lambda\}|}{\|F_T\|_1}
 \ge\left(1-\frac RT\right)\frac{\lambda|K|}{\|f\|_1}.
\]
令 $T\to\infty$，再令 $b\uparrow C_n$ 即得结论。此论证也给出有限维构造的具体承接方式；单调性见\cite{Aldaz2011}。
\end{proof}

\subsection{准确有限判据}
下界计算中的所有端点以整数及有理数表示。若 $m\in\N$、$D\in\N$，整数 $q$ 满足
\begin{equation}\label{pre:root}
 q^8\le mD^8<(q+1)^8,
\end{equation}
则 $q/D\le m^{1/8}<(q+1)/D$。和、差、积均按区间的包含方向传播；最后只用有理数严格比较。对见证盒的并集，先沿一个固定坐标的全部端点切片；每一开片上的截面盒集合不变。片宽乘截面并集体积后求和，再递归到一维区间合并，由 Fubini 定理与维数归纳得到准确体积。重复截面以计数维护，不将重叠盒的体积直接相加。

\section{主要证明}\subsection{几何覆盖与上界}\label{up:section}

弱型估计需要同时控制水平集的体积和所选平均方体的重叠。为兼顾这两点，我们按半径递减选择方体，以其新增体积定义权重，再用中心处的加权重叠决定后续选择。中心未被这一条件控制的方体，其新增部分位于一个方体角帽区域；该区域的加权体积可以通过坐标面的分解求出。这里采用的选择与覆盖思路与中心极大算子的经典几何方法有关，参见 \cite{SteinStromberg1983,AldazKissing2020}。下面给出所需选择过程及角帽估计的完整证明。

\subsubsection{参数与上界公式}

本节的底测度均为 $\R^n$ 上的 Lebesgue 测度。记
\[
 Q(y,r)=y+(-r,r)^n,\qquad
 M_nf(x)=\sup_{r>0}\frac{1}{(2r)^n}
             \int_{Q(x,r)}|f(y)|\dd y,
\]
其中各坐标使用同一个半径 $r$，并令
\[
 C_n=\sup_{\substack{f\in L^1(\R^n),\ \|f\|_1>0\\\lambda>0}}
 \frac{\lambda|\{M_nf>\lambda\}|}{\|f\|_1}.
\]
对 $t>1$，置
\begin{equation}\label{up:B}
\begin{split}
 q&=1+\frac1t,\qquad
 g_t(\rho)=1+\frac{1}{2t}-\frac{1}{2\rho},\\
 B_n(t)&=1-\left(\frac q2\right)^n+n\log t
       -\frac n2\int_1^t g_t(\rho)^{n-1}\frac{\dd\rho}{\rho}.
\end{split}
\end{equation}

\begin{theorem}\label{up:main}
对每个整数 $n\ge1$ 和每个 $t>1$，均有 $B_n(t)>0$，且
\begin{equation}\label{up:bound}
 C_n\le F_n(t):=\left[q^{n/2}+\sqrt{B_n(t)}\right]^2.
\end{equation}
特别地，取
\begin{equation}\label{up:t8}
 t_0=\frac{3837029}{524288},\qquad
 U=\frac{56627260639}{2000000000},
\end{equation}
则
\begin{equation}\label{up:eight}
 C_8\le F_8(t_0)<U=28.3136303195.
\end{equation}
\end{theorem}

证明时先保留一个有限的放大参数 $a>2t$。这样，覆盖中的每个权重和每个平均值都有实际对应的方体，待得到完整算子的估计后再令 $a$ 趋于无穷。固定 $\tau>0$，记
\begin{equation}\label{up:finite-parameters}
 b=aq,\qquad \alpha=\frac a{a+1},\qquad
 c=\frac{b-1}{a+1}>1,
\end{equation}
以及
\begin{equation}\label{up:finite-constants}
\begin{split}
 I_n(t)&=\int_1^t g_t(\rho)^{n-1}\frac{\dd\rho}{\rho},\\
 H_n(a,t)&=1-\left(\frac{\alpha q}{2}\right)^n
           +n\log t-\frac n2\alpha^{n-1}I_n(t),\\
 R_n(a,t)&=a^{-n}+\alpha^{-n}H_n(a,t),\qquad
 A_n(a,t)=\left(q+\frac1a\right)^n.
\end{split}
\end{equation}
我们将先证明
\begin{equation}\label{up:finite-operator}
 C_n\le\left(1+\frac{A_n(a,t)}{\tau}\right)
               \bigl(\tau+R_n(a,t)\bigr).
\end{equation}

\subsubsection{有限方体族的选择与覆盖}

若 $d>0$，以 $dQ$ 表示与 $Q$ 同心、半径扩大 $d$ 倍的方体。
给定任意有限个开方体，按如下过程选取其中的一部分。已有选取
$Q_j=Q(y_j,r_j)$ 时，固定其新增部分及权重
\begin{equation}\label{up:increments}
 D_j=Q_j\setminus\bigcup_{i<j}Q_i,\qquad
 v_j=|D_j|,\qquad w_j=\frac{v_j}{(2ar_j)^n}.
\end{equation}
在每一步，尚未选取且中心为 $y$ 的方体可选，当且仅当
\begin{equation}\label{up:eligible}
 \sum_{j\text{ 已选}}w_j\ind_{bQ_j}(y)\le\tau.
\end{equation}
每次从可选方体中选一个半径最大的；没有可选方体时停止。
每次选取后便移除一个待选方体，故过程在有限步内终止。权重在对应方体选入后保持不变。由于权重非负，一个方体一旦不可选便不会重新变为可选。因此，所选半径非增；新增集合 $D_j$ 两两不交，并满足
\begin{equation}\label{up:selected-volume}
 \sum_jv_j=\left|\bigcup_{j\text{ 已选}}Q_j\right|.
\end{equation}
这个过程允许权重为零，也允许半径相等。

\begin{lemma}\label{up:cover}
上述选择满足
\begin{equation}\label{up:cover-estimate}
 \left|\bigcup_{Q\text{ 给定}}Q\right|
 \le\left(1+\frac{A_n(a,t)}{\tau}\right)\sum_jv_j.
\end{equation}
\end{lemma}

\begin{proof}
设 $Q_\beta=Q(y_\beta,r_\beta)$ 最终未被选取。我们先说明
\begin{equation}\label{up:rejection}
 \sum_{j:r_j\ge r_\beta}w_j\ind_{bQ_j}(y_\beta)>\tau.
\end{equation}
若某个已选半径小于 $r_\beta$，考察第一次出现这种半径的时刻。
紧接该次选择之前，$Q_\beta$ 仍未被选入；若它可选，则按最大半径规则不能选择那个较小的方体。因此那时它已不可选，而此前所选半径均至少为 $r_\beta$，得到式 \eqref{up:rejection}。若从未选过半径小于 $r_\beta$ 的方体，则终止时的不可选条件直接给出同一结论。

对任意 $x\in Q_\beta$，式 \eqref{up:rejection} 中每个非零项都满足
\[
 \|x-y_j\|_\infty<r_\beta+br_j\le(b+1)r_j.
\]
故函数 $\sum_jw_j\ind_{(b+1)Q_j}$ 在未选方体的并上大于 $\tau$。
积分并加上已选方体的并，利用式 \eqref{up:selected-volume}，得
\[
 \left|\bigcup_{Q\text{ 给定}}Q\right|
 \le\sum_jv_j+\frac1\tau\sum_jw_j|(b+1)Q_j|
 =\left[1+\frac{((b+1)/a)^n}{\tau}\right]\sum_jv_j.
\]
由于 $(b+1)/a=q+1/a$，即得结论。这里扩大到 $b+1$ 倍包含了未选方体本身的半径。
\end{proof}

\subsubsection{重叠函数与未被中心条件控制的部分}

考虑所选方体的加权重叠函数
\begin{equation}\label{up:W}
 W(x)=\sum_jw_j\ind_{aQ_j}(x).
\end{equation}
我们将证明对每个 $x\in\R^n$ 均有
\begin{equation}\label{up:overlap}
 W(x)\le\tau+R_n(a,t).
\end{equation}
当 $W(x)>0$ 时，从对该值有正贡献的方体中选半径最小的一个；若有并列，取其中最后选入的下标 $k$。将 $x$ 平移到原点，并将所有坐标和半径同除以 $r_k$。Lebesgue 体积与分母同乘 $r_k^{-n}$，因而权重保持不变。于是可以假设 $x=0$、$r_k=1$。

其余有正贡献的项均满足 $r_j\ge1$ 且 $j<k$。事实上，若 $j>k$，半径非增给出 $r_j\le1$；此时只能等号成立，与并列时取最后下标的规则矛盾。方体 $Q_k$ 自身的权重满足
\begin{equation}\label{up:self}
 w_k\le\frac{|Q_k|}{|aQ_k|}=a^{-n}.
\end{equation}
对于其余正项，若 $y_k\in bQ_j$，它就在选择 $Q_k$ 时的非负和式 \eqref{up:eligible} 中，因此所有这类项的和不超过 $\tau$。

余下的下标满足
\begin{equation}\label{up:uncontrolled}
 \|y_k-y_j\|_\infty\ge br_j,\qquad
 \|y_j\|_\infty<ar_j,\qquad \|y_k\|_\infty<a.
\end{equation}
若 $r_j\ge t$，则
\[
 \|y_k-y_j\|_\infty<a+ar_j\le a(1+1/t)r_j=br_j,
\]
产生矛盾。因此这些下标的半径均在 $[1,t)$ 中，包含可能较早选入且半径等于 $1$ 的方体。

对这样的下标及 $y\in D_j\subset Q_j$，反三角不等式给出
\[
 \|y_k-y\|_\infty>(b-1)r_j.
\]
作变换
\[
 z=\frac y{a+1},\qquad u=\frac{y_k}{a+1},\qquad
 m(z)=\max(1,\|z\|_\infty).
\]
由 $0\in aQ_j$ 和 $1\le r_j<t$ 得
\[
 \|u\|_\infty<\alpha,\quad \|z\|_\infty<r_j<t,
 \quad m(z)\le r_j,\quad \|u-z\|_\infty>cr_j\ge cm(z).
\]
因而经变换后的新增集合均包含于
\begin{equation}\label{up:cap-domain}
 \Omega(u)=\{z:\|z\|_\infty<t,\ \|u-z\|_\infty\ge cm(z)\}.
\end{equation}
记其加权体积为
\begin{equation}\label{up:Psi}
 \Psi(u)=2^{-n}\int_{\Omega(u)}m(z)^{-n}\dd z.
\end{equation}
变换后的新增集合仍两两不交，而 $r_j^{-n}\le m(z)^{-n}$。利用 Lebesgue 变换的 Jacobian $(a+1)^n$，得到
\begin{equation}\label{up:cap-comparison}
 \sum_{j\text{ 满足 }\eqref{up:uncontrolled}}w_j
 \le\frac{(a+1)^n}{2^na^n}
          \int_{\Omega(u)}m(z)^{-n}\dd z
 =\alpha^{-n}\Psi(u).
\end{equation}
于是
\begin{equation}\label{up:W-Psi}
 W(0)\le\tau+a^{-n}+\alpha^{-n}\Psi(u).
\end{equation}
该比较使用的是新增集合的互不相交性；放大方体本身可以重叠。

\subsubsection{角点比较与方体壳层积分}

下面证明 $\Psi(u)$ 在闭方体 $\|u\|_\infty\le\alpha$ 上不超过角点处的值。坐标反射保持 Lebesgue 测度及 $m(z)$ 不变，故只需考虑
$h_i=|u_i|\in[0,\alpha]$。对 $\rho\ge1$，令
\[
 L_i(\rho)=\bigl|[-\rho,\rho]\cap[h_i-c\rho,h_i+c\rho]\bigr|.
\]
由于 $c>1$、$h_i\le\alpha<1\le\rho$，有
\begin{equation}\label{up:Li}
 L_i(\rho)=\min\{2\rho,(1+c)\rho-h_i\}
           =2\rho-[h_i-(c-1)\rho]_+,
\end{equation}
其中 $[s]_+=\max(s,0)$。它为正，且关于 $h_i$ 非增。
在 $\|z\|_\infty<1$ 内，$m(z)=1$，而 $\Omega(u)$ 的补集为两个开方体的交，故这一部分对 $\Psi$ 的贡献为
\begin{equation}\label{up:inner-cap}
 1-2^{-n}\prod_{i=1}^nL_i(1).
\end{equation}
此量关于每个 $h_i$ 非减。

在壳层 $\|z\|_\infty=\rho>1$ 上，按取得最大绝对值的坐标分解。
除去最大绝对值并列的零测集，每个点唯一属于一个坐标面；写该坐标为 $\rho$ 或 $-\rho$，其余坐标在 $(-\rho,\rho)$ 内自由变化。这一体积参数化的 Jacobian 绝对值为 $1$，完整壳层的面体积和为 $2n(2\rho)^{n-1}$。

我们称满足 $\|z-u\|_\infty<c\rho$ 的部分为壳层内的交叠部分。在第 $i$ 个正坐标面上，第 $i$ 个坐标始终满足此条件；负坐标面满足此条件当且仅当 $h_i<(c-1)\rho$。因此交叠部分的面体积和为
\begin{equation}\label{up:face-sum}
 S(\rho,h)=\sum_{i=1}^n
 \left(1+\ind_{\{h_i<(c-1)\rho\}}\right)
 \prod_{\ell\ne i}L_\ell(\rho).
\end{equation}
固定某个 $h_j$，每一项对它的依赖都只通过一个非负、非增因子：当 $i=j$ 时是括号内的面数，当 $i\ne j$ 时是 $L_j$。所以 $S$ 关于每个 $h_j$ 非增。

由此得到完整分解
\begin{equation}\label{up:shell}
\begin{split}
 \Psi(u)={}&1-2^{-n}\prod_{i=1}^nL_i(1)\\
 &+2^{-n}\int_1^t\rho^{-n}
  \bigl[2n(2\rho)^{n-1}-S(\rho,h)\bigr]\dd\rho.
\end{split}
\end{equation}
坐标面的边界不会改变面体积；等式 $h_i=(c-1)\rho$ 对每个坐标至多对应一个半径，也不会改变积分。该公式在 $n=1$ 时同样成立，约定空乘积为 $1$。
式 \eqref{up:inner-cap} 与式 \eqref{up:shell} 中的被积部分均关于每个 $h_j$ 非减，因而
\begin{equation}\label{up:corner}
 \Psi(u)\le\Psi(\alpha,\ldots,\alpha).
\end{equation}
实际中心对应 $\|u\|_\infty<\alpha$，比较角点可以位于闭方体边界。

接着求角点值。由参数关系
\[
 c-1=\frac{a/t-2}{a+1}>0,\qquad
 \frac{\alpha}{c-1}=\frac{at}{a-2t}>t,
\]
可知对 $1\le\rho\le t$，角点处每个坐标恰有一个坐标面属于交叠部分。此时
\[
 L(\rho)=(1+c)\rho-\alpha,
 \qquad S(\rho,\alpha\boldsymbol{1})=nL(\rho)^{n-1}.
\]
又因 $1+c=\alpha(2+1/t)$，有
\[
 \frac{L(\rho)}{2\rho}=\alpha g_t(\rho),\qquad
 \frac{L(1)}2=\frac{\alpha q}2.
\]
特别地，归一化后的交叠壳层贡献为
\[
 2^{-n}\rho^{-n}nL(\rho)^{n-1}\dd\rho
 =\frac n2\,[\alpha g_t(\rho)]^{n-1}\frac{\dd\rho}{\rho}.
\]
这里的系数 $n/2$ 来自每个坐标只有一个交叠面，而完整壳层每个坐标有两个面。代入式 \eqref{up:shell}，即得
\begin{equation}\label{up:corner-value}
 \Psi(\alpha,\ldots,\alpha)=H_n(a,t).
\end{equation}
由于 $q/2<1$，且 $0<g_t(\rho)\le1$，$H_n(a,t)$ 的内部贡献为正，其外部被积函数满足
\[
 \frac n\rho\left[1-\frac12(\alpha g_t(\rho))^{n-1}\right]
 \ge\frac n{2\rho}>0.
\]
结合式 \eqref{up:W-Psi}、\eqref{up:corner} 和 \eqref{up:corner-value}，证明了式 \eqref{up:overlap}。其中自身项 $a^{-n}$ 在所有有限 $a$ 的估计中均被保留。

\subsubsection{一般可积输入与全部正尺度}

现在将有限覆盖用于实际水平集。任取 $f\in L^1(\R^n)$ 和 $\lambda>0$，令
$E_\lambda=\{M_nf>\lambda\}$。固定 $R>0$ 时，中心平均关于中心连续，因为平移差满足
\[
 \left|\frac1{(2R)^n}\int_{Q(x+h,R)}|f|
       -\frac1{(2R)^n}\int_{Q(x,R)}|f|\right|
 \le(2R)^{-n}\bigl\||f|(\cdot+h)-|f|\bigr\|_1
 \longrightarrow0.
\]
这里使用了 $L^1$ 平移连续性。所有 $R>0$ 上的上确界因而是下半连续函数，故 $E_\lambda$ 为开集。

取任意紧集 $K\subset E_\lambda$。对每个 $x\in K$，选取实际半径 $R_x>0$，使 $Q(x,R_x)$ 上的平均大于 $\lambda$，并以
$Q(x,R_x/a)$ 作为待选小方体。它们覆盖 $K$，故有有限子覆盖。对这个有限族使用上述选择，每个选中方体的 $a$ 倍扩大仍是原先平均大于 $\lambda$ 的方体。因此
\begin{equation}\label{up:input}
\begin{split}
 \lambda\sum_jv_j
 &\le\sum_j\frac{v_j}{|aQ_j|}\int_{aQ_j}|f(y)|\dd y\\
 &=\int_{\R^n}|f(y)|W(y)\dd y
 \le\bigl(\tau+R_n(a,t)\bigr)\|f\|_1.
\end{split}
\end{equation}
和式有限且各项非负。将式 \eqref{up:cover-estimate} 用于该子覆盖，再用式 \eqref{up:input}，得到
\[
 \lambda|K|\le
 \left(1+\frac{A_n(a,t)}\tau\right)
 \bigl(\tau+R_n(a,t)\bigr)\|f\|_1.
\]
对 $E_\lambda$ 内的紧集取上确界，由 Lebesgue 测度的内正则性，得到同一式子中的 $|K|$ 可以替换为 $|E_\lambda|$。这一步同时证明水平集体积有限，故无须预先假定它有限。再对输入和水平取上确界，即得式 \eqref{up:finite-operator}。

这个论证对任意实值或复值 $L^1$ 输入成立，所选 $R_x$ 可以是任何正数。有限子覆盖仅用于一次几何估计，最终极大算子仍包含全部正半径。开、闭方体的边界均为 Lebesgue 零测集，因而两种约定给出相同平均。

\subsubsection{标量极限与阈值的选择}

固定 $n\ge1$、$t>1$、$\tau>0$，令有限参数 $a>2t$ 趋于无穷。
由式 \eqref{up:finite-constants}，因 $I_n(t)$ 与 $a$ 无关，得到
\[
 A_n(a,t)\longrightarrow q^n,\qquad
 R_n(a,t)\longrightarrow B_n(t).
\]
每个有限 $a$ 已经给出同一完整算子的估计 \eqref{up:finite-operator}，因此只需对其标量右端取极限，便有
\begin{equation}\label{up:scalar-limit}
 C_n\le\left(1+\frac{q^n}\tau\right)\bigl(\tau+B_n(t)\bigr).
\end{equation}
这里不要求随 $a$ 选出的子覆盖或权重收敛。

为确定阈值，先将式 \eqref{up:B} 改写为
\[
 B_n(t)=1-(q/2)^n+
 n\int_1^t\left[1-\frac12g_t(\rho)^{n-1}\right]\frac{\dd\rho}{\rho}>0.
\]
其首项为正，积分至少为 $(n/2)\log t$。令 $A=q^n$、$B=B_n(t)$，则式 \eqref{up:scalar-limit} 的右端为
\[
 A+B+\tau+\frac{AB}\tau.
\]
选取合法的正阈值 $\tau=\sqrt{AB}$，便得到 $(\sqrt A+\sqrt B)^2$，证明了式 \eqref{up:bound}。

\subsubsection{八维参数的准确比较}

八维常数的比较只涉及式 \eqref{up:bound} 中的标量函数。令
$p=1+1/(2t)$。二项式展开并逐项积分给出准确恒等式
\begin{equation}\label{up:scalar-expansion}
 I_8(t)=p^7\log t+
 \sum_{j=1}^7\binom7j p^{7-j}\left(-\frac12\right)^j
              \frac{1-t^{-j}}j.
\end{equation}
因此，当 $t$ 为有理数时，$B_8(t)$ 是一个有理数与
$(8-4p^7)\log t$ 的和。

在 $t=t_0$ 处，置 $z=(t_0-1)/(t_0+1)$，使用有理上下界
\begin{equation}\label{up:log-certificate}
 \ell=2\sum_{k=0}^{399}\frac{z^{2k+1}}{2k+1},\qquad
 \ell\le\log t_0\le
 \ell+\frac{2z^{801}}{801(1-z^2)}.
\end{equation}
此余项界由正项级数及几何级数直接得到：将余项中每个分母以 $801$ 从下方替换即可。将式 \eqref{up:log-certificate} 代入式 \eqref{up:scalar-expansion}，并按 $8-4p^7$ 的符号取相应端点，可得 $B_8(t_0)$ 的有理包围。其端点向外保留四十位小数时为
\begin{equation}\label{up:B-enclosure}
\begin{gathered}
 b_-<B_8(t_0)<b_+,\\
 b_-=13.3365553664751128154150225918012012532221,\\
 b_+=13.3365553664751128154150225918012012532222.
\end{gathered}
\end{equation}

严格平方根比较也可完全用有理数完成。将由式 \eqref{up:log-certificate} 得到的包围上端点向外保留四十位小数，记所得有理数为 $B_+$，并记
$s=(1+1/t_0)^4$、$d=U-s^2-B_+$。准确有理运算给出
\begin{equation}\label{up:residual}
 d>0,\qquad
 d^2-4s^2B_+>
 1.900403275813764719335\times10^{-9}>0.
\end{equation}
所用端点和整数比较可见附录的准确计算材料。由于 $s>0$、$B_+>0$，式 \eqref{up:residual} 蕴含
$d>2s\sqrt{B_+}$，故
\[
 F_8(t_0)\le(s+\sqrt{B_+})^2<U.
\]
这证明了式 \eqref{up:eight}，并完成定理 \ref{up:main} 的证明。

\begin{remark}\label{up:Lebesgue}
本证明中常 Lebesgue 底测度具体用于方体体积 $(2r)^n$、新增部分的体积权重、共同平移与伸缩下的归一化，以及坐标反射、交叠区域的乘积体积和壳层的体积参数化。一般 $L^1$ 输入的传递还使用了平移连续性、方体边界的零测性与内正则性。输入 $|f|$ 可以任意变化；这些几何体积公式依赖的是底测度的常密度。
\end{remark}
\subsection{格点构造与计算辅助下界}
\label{lo:section}

本节用一个有限的非负光滑函数证明下界。方法的出发点是
Aldaz~\cite{Aldaz2000} 的格点构造：在周期格点附近放置相同质量，
寻找能由同一中心立方体包含若干质量块的中心，再估计这些中心的总体积。
这里选取八个坐标方向上的两种间距，并把中心集合写成有限个轴平行盒的并。
盒的交叠通过准确的切片递归处理；最后直接构造有限光滑输入，控制平滑及截断造成的损失。

\begin{proposition}\label{lo:bound}
八维中心立方体极大算子的弱型常数满足
\[
 C_8>\frac{2188811179426669}{10^{15}}.
\]
这一严格下界可由一个非负的 $C_c^\infty(\R^8)$ 函数及一个明确的正水平值实现。
\end{proposition}

\paragraph{格点、源块与共同半径.}
令
\begin{equation}\label{lo:spacing}
 a_1=\frac{1+1944^{1/8}}2,\qquad
 a_2=\cdots=a_8=a=\frac{1+2^{1/8}}2,
 \qquad
 \Lambda=\{(a_1z_1,\ldots,a_8z_8):z\in\Z^8\}.
\end{equation}
记 $a\odot z=(a_1z_1,\ldots,a_8z_8)$，并令
\[
 Q=\prod_{i=1}^8[-a_i/2,a_i/2],\qquad
 Q_+=\prod_{i=1}^8[0,a_i/2].
\]
各个 $a\odot z+Q$ 构成 $\R^8$ 的基本胞分割，除边界外两两不交。
考虑长度标签集合
\[
 \mathcal E=\{1,2\}\times\{1,2,3,4\}^7,
 \qquad \#\mathcal E=32768.
\]
对 $\ell\in\mathcal E$，定义
\begin{equation}\label{lo:source-indices}
 q_i=\begin{cases}
       -\lfloor\ell_i/2\rfloor,&\ell_i\text{ 为奇数},\\
       1-\ell_i/2,&\ell_i\text{ 为偶数},
      \end{cases}
 \qquad p_i=q_i+\ell_i-1,
 \qquad
 J_\ell=\prod_{i=1}^8\{q_i,q_i+1,\ldots,p_i\}.
\end{equation}
这些是在正半胞方向使用的最近连续整数指标；总有 $q_i\le0\le p_i$。
源块 $\{a\odot z:z\in J_\ell\}$ 含有
$m_\ell=\prod_i\ell_i$ 个格点。取共同半径
\begin{equation}\label{lo:common-radius}
 r_\ell=\frac{m_\ell^{1/8}}2,\qquad (2r_\ell)^8=m_\ell.
\end{equation}
若 $x\in Q_+$，以 $x$ 为中心、半径为 $r_\ell$ 的立方体包含这个源块，当且仅当
\[
 p_i a_i-r_\ell\le x_i\le q_i a_i+r_\ell
 \quad(1\le i\le8).
\]
因而相应的中心盒为
\begin{equation}\label{lo:true-box}
 B_\ell=\prod_{i=1}^8
 \left[\max\{0,p_i a_i-r_\ell\},\,
       \min\{a_i/2,q_i a_i+r_\ell\}\right].
\end{equation}
只保留各边长均为正的盒。这里每个标签使用一个共同半径，
八个坐标条件同时保证整个源块被包含。

对 $\sigma\in\{-1,1\}^8$，记
$\sigma B=\{(\sigma_1x_1,\ldots,\sigma_8x_8):x\in B\}$。
格点及源块可按同一符号反射，故一个基本胞中的中心集合是
\[
 E=\bigcup_{\sigma\in\{-1,1\}^8}\ \bigcup_{\ell\in\mathcal E}\sigma B_\ell.
\]
不同正交象限的内部不交，因此
$|E|=2^8|\bigcup_\ell B_\ell|$。
下面准确估计这个有限并集的体积。

\paragraph{有理内盒和外盒.}
所有根式只需作有理包围。取 $D=10^{35}$、$S=4D$。
对每个出现的质量 $m$，令 $t_m$ 是满足
\begin{equation}\label{lo:root-certificate}
 t_m^8\le mD^8<(t_m+1)^8
\end{equation}
的整数。这两个整数幂不等式就是根式包围的检查判据。
首轴间距和其余间距的整数端点分别为
\begin{align*}
 A_1^-&=2(D+t_{1944}),& A_1^+&=2(D+t_{1944}+1),\\
 A_i^-&=2(D+t_2),& A_i^+&=2(D+t_2+1)\quad(2\le i\le8).
\end{align*}
于是 $A_i^-/S\le a_i\le A_i^+/S$。
首轴公式中的 $1^{1/8}=1$ 使用准确值，不再给这一项增加包围宽度。
对标签 $\ell$，令 $R_\ell^-=2t_{m_\ell}$、
$R_\ell^+=2(t_{m_\ell}+1)$，则
$R_\ell^-/S\le r_\ell\le R_\ell^+/S$。
即使质量是完全八次幂，这个半径包围仍然有效。

定义整数端点
\begin{align}
 L_{\ell i}^-&=\max\{0,p_iA_i^+-R_\ell^-\},&
 U_{\ell i}^-&=\min\{A_i^-/2,q_iA_i^++R_\ell^-\},\label{lo:inner-endpoints}\\
 L_{\ell i}^+&=\max\{0,p_iA_i^--R_\ell^+\},&
 U_{\ell i}^+&=\min\{A_i^+/2,q_iA_i^-+R_\ell^+\}.
 \label{lo:outer-endpoints}
\end{align}
$A_i^\pm$ 均为偶数，所以以上全为整数。令
\[
 B_\ell^\pm=\prod_{i=1}^8[L_{\ell i}^\pm/S,U_{\ell i}^\pm/S],
 \qquad
 E^\pm=\bigcup_{\sigma\in\{-1,1\}^8}\ \bigcup_\ell\sigma B_\ell^\pm,
\]
仍只保留正边长盒。由 $q_i\le0\le p_i$ 可逐项得到
\begin{equation}\label{lo:box-sandwich}
 B_\ell^-\subseteq B_\ell\subseteq B_\ell^+,
 \qquad |E^-|\le |E|\le |E^+|.
\end{equation}
外盒用于同时包围几何体积；实际输入的下界只使用内盒。

\begin{lemma}\label{lo:union-recursion}
有限个有理轴平行盒的并体积可由有限次有理数加法和乘法准确求出。
具体地，设 $\mathcal B$ 是 $d$ 维盒族，第一轴所有端点按递增次序排列为
$t_1<\cdots<t_N$。令 $\mathcal B_j$ 为在开条带
$(t_j,t_{j+1})\times\R^{d-1}$ 中存在的盒在后 $d-1$ 个坐标上的投影族。
若 $\mathcal U_d(\mathcal B)=|\bigcup_{B\in\mathcal B}B|$，则
\begin{equation}\label{lo:sweep-recursion}
 \mathcal U_d(\mathcal B)=
 \sum_{j=1}^{N-1}(t_{j+1}-t_j)\mathcal U_{d-1}(\mathcal B_j).
\end{equation}
在 $d=1$ 时，$\mathcal U_1$ 由合并有限个区间求出，空族的体积为零。
\end{lemma}

\begin{proof}
对 $x_1\in(t_j,t_{j+1})$，并集的截面恰为
$\bigcup_{B'\in\mathcal B_j}B'$，与 $x_1$ 无关。
由 Fubini 定理即得 \eqref{lo:sweep-recursion}；端点截面是零测集，不影响体积。
按维数归纳，递归终止于一维区间并。
所有端点为有理数，故每一步及最终结果都可用准确有理算术完成。
\end{proof}

实际计算先在整数坐标中使用上述递归，最后除以 $S^8$ 并乘以 $2^8$。
每层沿固定坐标次序按端点加入或移去盒。
若不同盒产生相同的后缀投影，须保留其重数；仅当重数降至零时才将该投影移去。
完全相同的物理盒可以预先去重，递归中完全相同的剩余盒集合可以缓存。
这些处理不改变引理中的截面族，因而不改变体积。

\begin{lemma}\label{lo:volume-certificate}
对 \eqref{lo:spacing}--\eqref{lo:outer-endpoints} 的准确输入，
有限并集计算给出
\begin{equation}\label{lo:compact-volume}
 \begin{gathered}
 L_0\le |E^-|\le |E|\le |E^+|\le U_0,\\
 L_0=\frac{21888111794266698273227523980123024}{10^{34}},\\
 U_0=\frac{21888111794266698273227523980123027}{10^{34}}.
 \end{gathered}
\end{equation}
内盒和外盒各有 $882$ 个正体积标签盒，去除完全相同的盒后各有 $644$ 个。
内盒的每条边长均大于 $2\cdot10^{-25}$，且小于 $2$。
\end{lemma}

\begin{proof}[计算辅助验证的判据与结果]
依次遍历 $\mathcal E$，由 \eqref{lo:source-indices} 计算整数指标与质量，
用 \eqref{lo:root-certificate} 检查根式包围，再按
\eqref{lo:inner-endpoints}、\eqref{lo:outer-endpoints} 生成全部盒。
以八个严格整数比较 $L_{\ell i}^\pm<U_{\ell i}^\pm$ 判定正体积，
并以八个端点对组成的元组识别完全相同的物理盒。
对两族盒分别应用引理~\ref{lo:union-recursion}，得到准确分数后作向外有理舍入，
得到 \eqref{lo:compact-volume}。
两次递归分别有 $104294$ 和 $106152$ 个不同状态。
全部 $882$ 个内盒的最小边长为
\[
 \frac{269425567040611903127866148065111}
      {40000000000000000000000000000000000}
 >2\cdot10^{-25}.
\]
边长小于 $2$ 也由整数端点比较得到。
补充材料提供准确输入、完整体积分数及检查程序；
其中 \texttt{supplement/paper\_claims.json} 保存本文采用的准确数值。
\end{proof}

\paragraph{有限光滑输入与严格水平集.}
为把几何体积用于 $L^1$ 弱型比值，需要同时控制格点质量的平滑和格点族的有限截断。
取
\begin{equation}\label{lo:finite-parameters}
 \varepsilon=10^{-25},\qquad \eta=\varepsilon/2,\qquad
 \theta=1-10^{-18},\qquad K=10^{19}.
\end{equation}
对每个正体积内盒，把每一面的坐标向内移动 $\varepsilon$，记所得盒为
\[
 B_{\ell,\varepsilon}^-
 =\prod_{i=1}^8
 [L_{\ell i}^-/S+\varepsilon,U_{\ell i}^-/S-\varepsilon],
 \qquad
 E_\varepsilon^-=\bigcup_\sigma\ \bigcup_\ell\sigma B_{\ell,\varepsilon}^-.
\]
引理~\ref{lo:volume-certificate} 保证这些盒仍有正边长。

\begin{lemma}\label{lo:smoothing-loss}
有
\begin{equation}\label{lo:loss}
 |E_\varepsilon^-|\ge L_0-\Delta,\qquad
 \Delta=2^8\cdot882\cdot16\varepsilon\cdot2^7
       =\frac{441}{9536743164062500000}.
\end{equation}
若 $\varphi_\eta\ge0$ 是积分为 $1$、支撑于 $[-\eta,\eta]^8$ 的光滑函数，
则对 $x\in B_{\ell,\varepsilon}^-$，同一个立方体
$x+[-r_\ell,r_\ell]^8$ 包含全部
$\varphi_\eta(\,\cdot-a\odot z)$（$z\in J_\ell$）的支撑。
相应结论在每个反射盒上仍成立。
\end{lemma}

\begin{proof}
每个内盒有 $16$ 个面，且各边长小于 $2$。
向内移动一个面造成的损失包含于体积至多 $\varepsilon2^7$ 的薄片中。
先对一个盒的各面求和，再对 $882$ 个标签及 $2^8$ 个反射求和，
由测度的次可加性得到 \eqref{lo:loss}。

对 $x\in B_{\ell,\varepsilon}^-$，
\eqref{lo:box-sandwich} 及各面的内缩给出
\[
 p_i a_i-r_\ell+\varepsilon\le x_i
 \le q_i a_i+r_\ell-\varepsilon.
\]
故每个 $z\in J_\ell$ 满足
$|x_i-a_i z_i|\le r_\ell-\varepsilon$。
支撑中各点与相应格点的坐标距离不超过 $\eta<\varepsilon$，
于是整个支撑包含在同一个中心立方体内。
反射保持这些坐标距离，结论随之成立。
\end{proof}

\begin{proof}[命题~\ref{lo:bound} 的证明]
固定一个偶的非负函数 $\varphi_\eta\in C_c^\infty(\R^8)$，
使其积分为 $1$、支撑包含于 $[-\eta,\eta]^8$。
这里的偶性是对每个坐标分别而言。
令
\begin{equation}\label{lo:finite-input}
 F(x)=\sum_{z\in\{-K-2,\ldots,K+2\}^8}
       \varphi_\eta(x-a\odot z).
\end{equation}
这是有限项之和，因而 $F\ge0$、$F\in C_c^\infty(\R^8)$，且
\begin{equation}\label{lo:input-mass}
 \|F\|_1=(2K+5)^8.
\end{equation}
由 \eqref{lo:source-indices}，所有源指标属于 $[-1,2]^8\cap\Z^8$；
反射后的指标属于 $[-2,2]^8\cap\Z^8$。
因此，对每个接收格点 $z^0\in\{-K,\ldots,K\}^8$，
各反射源块平移 $a\odot z^0$ 后仍全部包含于 \eqref{lo:finite-input} 的源指标范围中。
每个坐标增加两层格点已足够。

若 $x\in a\odot z^0+E_\varepsilon^-$，选取包含它的一个平移反射盒，
并使用对应的 $r_\ell$。
引理~\ref{lo:smoothing-loss} 保证该中心立方体包含 $m_\ell$ 个完整质量块，
由 \eqref{lo:common-radius} 可得
\[
 M_8F(x)\ge
 \frac{1}{(2r_\ell)^8}
 \int_{x+[-r_\ell,r_\ell]^8}F(y)\,\dd y
 \ge\frac{m_\ell}{(2r_\ell)^8}=1>\theta.
\]
不同接收格点的集合 $a\odot z^0+E_\varepsilon^-$ 位于不同基本胞内，
除零测边界外两两不交。因此
\[
 |\{x:M_8F(x)>\theta\}|
 \ge(2K+1)^8|E_\varepsilon^-|
 \ge(2K+1)^8(L_0-\Delta).
\]
将 \eqref{lo:input-mass} 代入弱型比值，得到
\begin{equation}\label{lo:finite-ratio}
 C_8\ge
 \frac{\theta|\{M_8F>\theta\}|}{\|F\|_1}
 \ge\theta(L_0-\Delta)
       \left(\frac{2K+1}{2K+5}\right)^8
 >\frac{2188811179426669}{10^{15}}.
\end{equation}
最后一个严格比较把 \eqref{lo:compact-volume}、\eqref{lo:finite-parameters}
和 \eqref{lo:loss} 中的有理数代入后，通分为整数不等式得到；
补充材料同时保存该比较的准确分数与检查程序。
全部见证半径均为正数，水平值严格小于 $1$，输入本身是有限光滑函数，
故这给出所定义的一般非负 $L^1(\R^8)$ 输入类中的下界。
\end{proof}
\subsection{有限方向势的奇异分层与卷积比较}
\label{gr:sec}

本节建立从有限方向 Cauchy 半群到立方体极大算子的比较原理。
方法取自 Lin 的势核与障碍分解思路\cite{Lin2026}；下面针对八维有限方向模型给出完整证明。
其要点是先证明 Green 势的正部具有全空间的有限正源，再用一个与半径无关的障碍集同时控制所有卷积尺度。
有限方向造成的低秩奇异集必须在这一过程中一并处理。

\subsubsection{半群、生成元与奇异集}

设 $v_1,\ldots,v_m\in\R^8\setminus\{0\}$ 张成 $\R^8$，$c_j>0$。
方向允许平行或重复，以下始终使用给定向量，不将其单位化。
记
\[
 q_a(s)=\frac{a}{\pi(a^2+s^2)},\qquad
 X_t=\sum_{j=1}^m v_jZ_{j,t},
 \qquad Z_{j,t}\sim q_{c_jt},
\]
其中各系数独立。令 $p_t$ 为 $X_t$ 的密度，并定义
\begin{equation}\label{gr:objects}
 \psi(\xi)=\sum_{j=1}^m c_j|v_j\cdot\xi|,
 \qquad P_tf=p_t*f,\qquad L=-\psi(D),
 \qquad G(x)=\int_0^\infty p_t(x)\,\dd t.
\end{equation}
本节采用 Fourier 变换
$\widehat f(\xi)=\int_{\R^8}e^{-ix\cdot\xi}f(x)\,\dd x$，
逆变换的系数为 $(2\pi)^{-8}$。所有积分、卷积及弱型估计均以 Lebesgue 测度为背景。

取任一由给定方向组成的基 $B=(v_{i_1},\ldots,v_{i_8})$，其部分和的密度为
\begin{equation}\label{gr:basis-density}
 p_t^B(x)=\frac1{|\det B|}
 \prod_{k=1}^8\frac{c_{i_k}t}
 {\pi\bigl((c_{i_k}t)^2+(B^{-1}x)_k^2\bigr)}.
	ag{\ref{gr:basis-density}}
\end{equation}
将其与其余项的概率律卷积，即得真实密度 $p_t$。因此 $p_t>0$，$p_t\in C_0(\R^8)$，$\int p_t=1$，且
\begin{equation}\label{gr:global-density}
 \|p_t\|_\infty\le Ct^{-8},\qquad
 C=\frac1{\pi^8|\det B|\prod_{k=1}^8c_{i_k}}.
\end{equation}
各系数的缩放还给出
\begin{equation}\label{gr:scaling}
 p_t(x)=t^{-8}p_1(x/t),\qquad G(rx)=r^{-7}G(x)\quad(r>0),
\end{equation}
其中第二式允许两边取无穷值。
概率构造表明 $(P_t)_{t\ge0}$ 是自伴、保正且保持质量的卷积半群；$p_t$ 是近似恒等核。
缩放及 $p_1\in C_0\cap L^1$ 又保证正时间的联合连续性及 $L^1$ 连续性。
由方向张成性及单位球的紧性，存在 $0<a\le b<\infty$ 使
\begin{equation}\label{gr:ellipticity}
 a|\xi|\le\psi(\xi)\le b|\xi|.
\end{equation}
因此 $L$ 的 $L^2$ 算子域为 $H^1$，形式域为 $H^{1/2}$；这里的范数均包含 $L^2$ 项。

对 $\phi\in C_c^\infty(\R^8)$，生成元的实空间表达为
\begin{equation}\label{gr:generator}
 L\phi(x)=\frac1{2\pi}\sum_{j=1}^m c_j
 \int_\R\frac{\phi(x+sv_j)+\phi(x-sv_j)-2\phi(x)}{s^2}\,\dd s.
\end{equation}
其 Fourier 符号恰为 $-\psi$；特别地，式中系数为 $c_j$，给定向量的长度保留在位移中。
二阶差分满足
\[
 \|\phi(\cdot+sv_j)+\phi(\cdot-sv_j)-2\phi\|_\infty
 \le\min\{s^2|v_j|^2\|D^2\phi\|_\infty,4\|\phi\|_\infty\},
\]
并有对应的 $L^1$ 估计，故 $L\phi\in L^1\cap L^\infty\cap C_0$。
一阶齐次性给出，对于 $\chi_R(x)=\chi(x/R)$，
\begin{equation}\label{gr:cutoff-scaling}
 L\chi_R(x)=R^{-1}(L\chi)(x/R),\quad
 \|L\chi_R\|_\infty=R^{-1}\|L\chi\|_\infty,\quad
 \|L\chi_\varepsilon\|_1=\varepsilon^7\|L\chi\|_1.
\end{equation}
Fourier 支配收敛还给出，对 Schwartz 函数 $\phi$，
\begin{equation}\label{gr:test-generator}
 \frac{P_t\phi-\phi}{t}\longrightarrow L\phi
 \quad\hbox{于 }L^\infty\hbox{ 及 }L^2,
 \qquad \partial_tP_t\phi=P_tL\phi.
\end{equation}
这些测试函数上的性质将用于处理全空间分布等式。

\begin{theorem}\label{gr:singular}
定义
\[
 S_3=\bigcup_{\substack{A\subseteq\{v_1,\ldots,v_m\}\\ \rank A\le3}}
 \spanop(A),
\]
其中包括空集张成的 $\{0\}$。则 $S_3$ 为闭的零测集，且
\[
 G(x)=+\infty\quad\Longleftrightarrow\quad x\in S_3.
\]
$G$ 在 $\R^8\setminus S_3$ 上有限且连续。
\end{theorem}

\begin{proof}
$S_3$ 是有限个维数至多为三的闭线性子空间之并，故闭且为零测集。
先取 $E\Subset\R^8\setminus S_3$，记
\[
 d=\operatorname{dist}(E,S_3)>0,\qquad
 D_V=\sum_{j=1}^m|v_j|,\qquad \delta=\frac d{4D_V},
 \qquad O=\{x:\operatorname{dist}(x,E)<d/4\}.
\]
对系数 $s=(s_j)$，置 $J(s)=\{j:|s_j|>\delta\}$。
若这些大系数对应的方向秩至多为三，则其和属于 $S_3$，而
\[
 \left|\sum_{j\notin J(s)}s_jv_j\right|\le\delta D_V=d/4.
\]
此时 $\sum_js_jv_j$ 到 $S_3$ 的距离至多为 $d/4$，不可能落入 $O$。
所以，一个落入 $O$ 的表示中，全部大系数方向的秩至少为四。

令 $\mathcal I_4$ 为独立四元方向指标集的全体，对每个 $I\in\mathcal I_4$ 选一个包含它的方向基 $B_I$。
将 $I$ 中四个系数的概率律限制在 $|s|>\delta$，保留其未经归一化的子概率律。
这些密度满足
\[
 q_{c_jt}(s)\ind_{\{|s|>\delta\}}\le\frac{c_jt}{\pi\delta^2}.
\]
基中其余四个系数的密度至多为 $(\pi c_jt)^{-1}$。
由基变换及与其余概率律卷积，受限随机和的密度至多为
\[
 C_I=\frac{\prod_{j\in I}c_j}
 {\pi^8|\det B_I|\delta^8\prod_{j\in B_I\setminus I}c_j}.
\]
这里 $t^4t^{-4}=1$；没有因条件化而除以限制事件的概率。
对任意 Borel 集 $F\subset O$，大系数的秩结论与事件并集给出
\[
 \mathbb P(X_t\in F)
 \le\sum_{I\in\mathcal I_4}
 \mathbb P\bigl(X_t\in F,\ |Z_{j,t}|>\delta\ (j\in I)\bigr)
 \le\Bigl(\sum_I C_I\Bigr)|F|.
\]
因此 $p_t\le C_E:=\sum_I C_I$ 在 $O$ 上几乎处处成立，真实密度的连续性将其提升为逐点界。
短时间用此界，长时间用 \eqref{gr:global-density}，便得到在 $E$ 上一致可积的时间控制函数。
支配收敛证明 $G$ 在 $E$ 上有限连续。

反之，若 $x\in S_3$，可写成
$x=\sum_{k=1}^ra_kv_{i_k}$，其中 $0\le r\le3$，且所列方向独立。
将它们补成方向基 $B$，则 $B^{-1}x=(a_1,\ldots,a_r,0,\ldots,0)$。
其余项之和记作 $Y_t$；由 $Y_t\overset d=tY_1$，可取 $M<\infty$，使
\[
 \mathbb P\{\|B^{-1}Y_t\|_\infty\le Mt\}=\theta>0
\]
与 $t$ 无关；没有其余项时取 $\theta=1$。
在这一事件上，$0<t\le1$ 时，\eqref{gr:basis-density} 的前 $r$ 个因子均有形如 $k_xt$ 的正下界，
因为相应分母由 $c_{i_k}^2+(|a_k|+M)^2$ 控制；后 $8-r$ 个因子均有形如 $k/t$ 的下界。
对 $Y_t$ 取期望，得到
\[
 p_t(x)=\mathbb E[p_t^B(x-Y_t)]\ge k'_xt^{2r-8},\qquad 0<t\le1.
\]
由于 $2r-8\le-2$，小时间积分发散，故 $G(x)=\infty$。
\end{proof}

\subsubsection{全空间可积正部及无原子概率源}

后面的卷积比较需要一个可积核，其生成元是正概率测度与原点质量之差。
我们先用半群估计获得可积性，再在全空间建立这一等式。

\begin{proposition}\label{gr:source}
令 $q=8/7$，$A_8=2(2C/7)^{1/7}$。则
\begin{equation}\label{gr:weak-green}
 |\{G>\lambda\}|\le A_8\lambda^{-8/7}\qquad(\lambda>0),
\end{equation}
且 $G\in L^1_{\mathrm{loc}}(\R^8)$。
对于每个 $c>0$，$H_c=(G-c)_+\in L^1(\R^8)$，并存在无原子的正概率测度 $\mu_c$，使
\begin{equation}\label{gr:source-identity}
 LH_c=\mu_c-\delta_0\qquad\hbox{于 }\R^8\hbox{ 的分布意义下成立}.
\end{equation}
\end{proposition}

\begin{proof}
写
\[
 U_T=\int_0^Tp_t\,\dd t,\qquad W_T=\int_T^\infty p_t\,\dd t.
\]
由质量与密度界，$\|U_T\|_1=T$，$\|W_T\|_\infty\le C/(7T^7)$。
取 $T=(2C/(7\lambda))^{1/7}$，则 $W_T\le\lambda/2$，Chebyshev 不等式给出 \eqref{gr:weak-green}。
对于有限测度集 $E$，Tonelli 定理给出
\[
 \int_EG\le\int_0^\infty\min\{1,C|E|t^{-8}\}\,\dd t<\infty.
\]
因此 $G$ 局部可积；层积分又给出
\begin{equation}\label{gr:positive-integrable}
 \int_{\R^8}H_c
 =\int_c^\infty|\{G>s\}|\,\dd s
 \le7A_8c^{-1/7}<\infty.
\end{equation}
这同时控制了奇异集邻域和空间无穷远。

对 $\phi\in C_c^\infty$，由 $L\phi\in L^1\cap L^\infty$，
\[
 \int G|L\phi|
 \le\int_0^\infty
 \min\{\|L\phi\|_\infty,Ct^{-8}\|L\phi\|_1\}\,\dd t<\infty.
\]
因而可以积分半群导数。自伴性及 $P_T\phi(0)\to0$ 给出
\begin{equation}\label{gr:fundamental}
 \langle G,L\phi\rangle
 =\int_0^\infty P_tL\phi(0)\,\dd t=-\phi(0),
 \qquad LG=-\delta_0.
\end{equation}
由 Tonelli 定理及半群律，
$P_tG=\int_t^\infty p_s\,\dd s$，故在 $G$ 几乎处处有限的点上 $G=P_tG+U_t$。
Jensen 不等式及正部函数的单调性、$1$-Lipschitz 性质给出
\[
 P_tH_c\ge(P_tG-c)_+\ge H_c-U_t.
\]
将此式与非负紧光滑测试函数配对，除以 $t$ 后令 $t\downarrow0$。
利用 \eqref{gr:test-generator}、$H_c\in L^1$ 及 $U_t/t\rightharpoonup\delta_0$，可得
$LH_c+\delta_0\ge0$。
正分布定理说明 $\mu_c:=LH_c+\delta_0$ 是全空间的正 Radon 测度。
该论证同时涵盖所有低秩子空间和水平界面。

选径向单调的 $\chi\in C_c^\infty$，$0\le\chi\le1$，在原点邻域等于一，使 $\chi_R\uparrow1$。
由 \eqref{gr:cutoff-scaling}，
\[
 \int\chi_R\,\dd\mu_c
 =1+\int H_cL\chi_R\longrightarrow1.
\]
正测度的单调收敛定理于是给出 $\mu_c(\R^8)=1$。
最后令 $h=\min(G,c)$。由已经证明的配对可积性及 $G=H_c+h$，有 $\mu_c=-Lh$。
对任意 $x_0\in\R^8$，使用平移后的小截断，得到
\[
 0\le\mu_c(\{x_0\})
 \le\int\chi\bigl((x-x_0)/\varepsilon\bigr)\,\dd\mu_c(x)
 =-\int h(x)L\bigl[\chi((\cdot-x_0)/\varepsilon)\bigr](x)\,\dd x
 \le c\varepsilon^7\|L\chi\|_1\longrightarrow0.
\]
故 $\mu_c$ 没有原子。
\end{proof}

\subsubsection{障碍分解与有限源卷积比较}

下面证明的分解只依赖于半群的质量守恒、对称性和前述密度界。
它产生一个固定集合，随后每个尺度的卷积比较都在该集合的补集上进行。

\begin{lemma}\label{gr:obstacle}
设 $f\ge0$ 有界且紧支撑，$\kappa>0$。存在
$u\ge0$，$u\in L^1\cap L^\infty\cap H^1$，及 $\sigma\in L^1$，使
\begin{equation}\label{gr:obstacle-result}
 f=\sigma-Lu,\qquad 0\le\sigma\le\kappa,\qquad
 \sigma=\kappa\ \hbox{几乎处处于 }D:=\{u>0\},\qquad
 \int\sigma=\int f.
\end{equation}
特别地，$\kappa|D|\le\|f\|_1$。
\end{lemma}

\begin{proof}
零输入的结论显然，以下令 $F=\|f\|_\infty>0$，$m_f=\|f\|_1>0$。
对于 $0<t\le1$，置
\[
 V_t=t\sum_{j\ge0}P_{jt}f,\qquad
 u_{t,0}=0,\qquad
 u_{t,k+1}=\bigl(P_tu_{t,k}+t(f-\kappa)\bigr)_+.
\]
由 $\|P_sf\|_\infty\le\min\{F,Cm_fs^{-8}\}$，正时间级数的尾项在上确界范数下收敛，且 $V_t-tf\in C_0$。
比较递减函数的右 Riemann 和与积分，得
\[
 \|V_t\|_\infty
 \le tF+\frac87(Cm_f)^{1/8}F^{7/8}\le C_f,
 \qquad C_f:=F+\frac87(Cm_f)^{1/8}F^{7/8}.
\]
Tonelli 定理给出 $P_tV_t=V_t-tf$。
迭代关于 $k$ 单调，且由 $V_t$ 控制，其极限满足
\[
 u_t=\bigl(P_tu_t+t(f-\kappa)\bigr)_+,
 \qquad 0\le u_t\le(V_t-\kappa t)_+.
\]
固定 $t$ 时，右侧有本质紧支撑，所以 $u_t\in L^1\cap L^2$。
这里没有要求 $V_t\in L^1$，也不需要不同 $t$ 的支撑包含在同一紧集中。

令 $L_t=(P_t-I)/t$，$\sigma_t=f+L_tu_t$。
固定点方程的正分支给出 $\sigma_t=\kappa$ 于 $\{u_t>0\}$；零分支给出
$0\le\sigma_t\le\kappa$ 于 $\{u_t=0\}$。
因 $u_t\in L^1$，质量守恒保证 $\int\sigma_t=m_f$。
将方程与 $u_t$ 配对可得
\[
 \langle u_t,-L_tu_t\rangle+\kappa\|u_t\|_1=\int fu_t.
\]
自伴收缩性使第一项非负，因而
\begin{equation}\label{gr:obstacle-bounds}
 \|u_t\|_1\le\frac{m_fC_f}{\kappa},\quad
 \|u_t\|_2^2\le\frac{m_fC_f^2}{\kappa},\quad
 \|\sigma_t\|_2^2\le\kappa m_f,\quad
 \|L_tu_t\|_2\le\|f\|_2+\sqrt{\kappa m_f}=:C_2.
\end{equation}

为获得局部强紧性，在本段采用酉 Fourier 归一化
$\mathcal F_uf=(2\pi)^{-4}\widehat f$，从而 Plancherel 等式不含额外系数。
记 $\psi_t=(1-e^{-t\psi})/t$。
若 $t\le1/Q$，则在 $\{\psi\ge Q\}$ 上
$\psi_t\ge(1-e^{-1})Q$，由 \eqref{gr:obstacle-bounds} 得
\begin{equation}\label{gr:high-frequency}
 \int_{\{\psi\ge Q\}}|\mathcal F_uu_t|^2
 \le\frac{C_2^2}{(1-e^{-1})^2Q^2}.
\end{equation}
固定低频投影由 \eqref{gr:ellipticity} 具有一致的 $H^1$ 界。
对其在有界区域使用 Rellich 紧嵌入，再用 \eqref{gr:high-frequency} 控制高频尾项，并作对角取子列，
可取 $t_j\downarrow0$ 使
\[
 u_{t_j}\to u\quad\hbox{于 }L^2_{\mathrm{loc}}\hbox{ 及几乎处处},
 \qquad u_{t_j}\rightharpoonup u,\quad
 \sigma_{t_j}\rightharpoonup\sigma\quad\hbox{于 }L^2.
\]
Fatou 引理及一致界保证 $u\ge0$，$u\in L^1\cap L^\infty$。
测试函数上 $L_t\phi\to L\phi$ 于 $L^2$，故极限满足 $Lu=\sigma-f\in L^2$。
由自伴 Fourier 乘子 $-\psi$ 的算子域刻画及 \eqref{gr:ellipticity}，得到 $u\in H^1$。
这里也可先利用 $C_c^\infty$ 在生成元图范数中的稠密性，再使用算子的闭性。

弱极限仍满足 $0\le\sigma\le\kappa$。
在 $D=\{u>0\}$ 的几乎处处点，$u_{t_j}$ 最终为正，故 $\sigma_{t_j}$ 最终等于 $\kappa$。
在 $D$ 的有限测度截集上用支配收敛，并与弱极限比较，得 $\sigma=\kappa$ 于 $D$。
由非负紧截断先得 $\sigma\in L^1$，$\int\sigma\le m_f$。
再以 $\chi_R$ 测试 $f=\sigma-Lu$；由
\[
 |\langle Lu,\chi_R\rangle|
 =|\langle u,L\chi_R\rangle|
 \le R^{-1}\|u\|_1\|L\chi\|_\infty\longrightarrow0,
\]
令 $R\to\infty$，得到 $\int\sigma=m_f$，从而排除了极限中的质量损失。
最后由 $\sigma=\kappa$ 于 $D$，有 $\kappa|D|\le m_f$。
\end{proof}

\begin{theorem}\label{gr:comparison}
设 $K\ge0$，$K\in L^1(\R^8)$，$M_K=\int K>0$，且对某个 $a>0$ 和正概率测度 $\mu$，有全空间分布等式
\begin{equation}\label{gr:finite-source}
 LK=a(\mu-\delta_0).
\end{equation}
记 $K_r(x)=r^{-8}K(x/r)$。则对任意 $f\in L^1(\R^8)$、$\lambda>0$，
\begin{equation}\label{gr:kernel-weak}
 \lambda\left|\left\{x:\sup_{r>0}K_r*|f|(x)>\lambda\right\}\right|
 \le M_K\|f\|_1.
\end{equation}
\end{theorem}

\begin{proof}
首先为卷积等式给出合法的 Fourier 解释。
对 Schwartz 测试 $\phi$，令 $\phi_R=\chi_R\phi$。
有 $\phi_R\to\phi$ 于 Schwartz 拓扑，二阶差分估计给出 $L\phi_R\to L\phi$ 于上确界范数。
由 $K\in L^1$ 及 $\mu$ 有限，\eqref{gr:finite-source} 因而延伸至 Schwartz 测试。
Fourier 反演和 Fubini 定理于是给出连续函数恒等式
\[
 -\psi(\xi)\widehat K(\xi)=a\bigl(\widehat\mu(\xi)-1\bigr).
\]
将此式乘以 $u\in H^1$ 的酉 Fourier 变换，利用 $L^1*L^2\to L^2$ 的 Young 不等式及有限测度卷积的 $L^2$ 界，得到
\begin{equation}\label{gr:convolution-identity}
 K*Lu=a(\mu*u-u)\qquad\hbox{于 }L^2.
\end{equation}
令 $\mu_r$ 为 $z\mapsto rz$ 对 $\mu$ 的推前测度。一阶齐次性给出精确缩放
\begin{equation}\label{gr:source-dilation}
 LK_r=\frac ar(\mu_r-\delta_0),\qquad
 K_r*Lu=\frac ar(\mu_r*u-u).
\end{equation}

先设 $f\ge0$ 有界紧支撑，以 $\kappa=\lambda/M_K$ 使用引理~\ref{gr:obstacle}。
对每个固定 $r>0$，在同一个 $D^c=\{u=0\}$ 上有 $\mu_r*u-u\ge0$ 几乎处处。
即使 $\mu_r$ 奇异，关于 $\dd x\,\dd\mu_r$ 的 Fubini 定理仍保证卷积对 $u$ 的零集修改不敏感，且保留其非负性。
由 \eqref{gr:source-dilation}，
\[
 K_r*f=K_r*\sigma-K_r*Lu
 \le K_r*\sigma\le\kappa M_K=\lambda
 \qquad\hbox{几乎处处于 }D^c.
\]
先对正有理数 $r$ 合并例外零集。
由于 $r\mapsto K_r$ 在 $L^1$ 中连续，且
\[
 \|K_r*f-K_s*f\|_\infty
 \le\|f\|_\infty\|K_r-K_s\|_1\longrightarrow0\quad(s\to r),
\]
上述结论延伸至全部正实半径。
因此严格超水平集模零集包含于同一个 $D$，从而
\[
 \lambda\bigl|\{\sup_{r>0}K_r*f>\lambda\}\bigr|
 \le\lambda|D|\le M_K\|f\|_1.
\]

一般输入取 $f_j=\min(|f|,j)\ind_{[-j,j]^8}\uparrow|f|$。
将卷积理解为非负的扩展值积分，由单调收敛及两个上确界的交换，
\[
 \sup_{r>0}K_r*f_j\uparrow\sup_{r>0}K_r*|f|,
 \qquad \sup_r\sup_jK_r*f_j=\sup_j\sup_rK_r*f_j.
\]
各有界近似输入的全半径上确界等于有理半径上确界，故可测，其单调极限亦可测。
对严格超水平集和输入质量应用单调收敛，便得到 \eqref{gr:kernel-weak}。
\end{proof}

\subsubsection{立方体比较与面上质量公式}

记 $Q=[-1,1]^8$，并令
\[
 J=\sum_{F\subset\partial Q}\int_FG(y)^{8/7}\,\dd S(y),
 \qquad g=\essinf_{x\in Q}G(x),
\]
其中和遍历十六个面；面的交界在表面积意义下可忽略。
在每个开面上，射线坐标 $x=ry$ 的 Jacobian 为 $r^7\,\dd r\,\dd S(y)$。
不同开面所生成的锥仅在零体积集上相交，所以由 Tonelli 定理及 \eqref{gr:scaling}，
\begin{equation}\label{gr:mass-formula}
 |\{G>1\}|=\frac18J,\qquad
 |\{G>s\}|=\frac18Js^{-8/7},\qquad
 \int H_c=\frac78Jc^{-1/7}.
\end{equation}
命题~\ref{gr:source} 保证第一式左侧有限，故 $J<\infty$；这里没有预先假定面上可积性。

$p_t$ 在 $Q\times[1,2]$ 上严格正且连续，因而
$\min_{x\in Q}\int_1^2p_t(x)\,\dd t>0$。
结合局部可积性可知 $0<g<\infty$。
并且
\begin{equation}\label{gr:essential-face}
 g=\essinf_{y\in\partial Q}G(y).
\end{equation}
事实上，若边界几乎处处有下界 $b$，则 $r^{-7}\ge1$ 立即给出内部同一下界。
反之，若边界上有正表面积的点集满足 $G(y)\le a<g$，取一个充分接近一的区间
$r_0<r<1$，使 $r_0^{-7}a<g$，则这些射线在 $Q$ 内生成正体积且 $G<g$ 的点集，矛盾。
由定理~\ref{gr:singular}，在当前有限方向情形，任何本质立方体下界还都是逐点下界：
$S_3$ 上势为无穷，而补集上的连续性排除了低于该下界的点。

\begin{corollary}\label{gr:cube}
设 $0<g_0\le g$。中心轴平行立方体极大算子
\[
 M_Qf(x)=\sup_{r>0}\frac1{(2r)^8}\int_{x+rQ}|f(y)|\,\dd y
\]
的弱 $(1,1)$ 常数满足
\begin{equation}\label{gr:conditional-J}
 C_8\le\frac{8^{1/7}J}{256g_0^{8/7}}.
\end{equation}
若模型在各坐标反射下不变，且其对称群在十六个面上传递，记
\[
 I=\int_{[0,1]^7}G(1,y)^{8/7}\,\dd y.
\]
则 $J=2048I$；对于任意上界 $I\le I_0$，有条件估计
\begin{equation}\label{gr:conditional-I}
 C_8\le\frac{8^{8/7}I_0}{g_0^{8/7}},\qquad
 C_8^7\le\frac{8^8I_0^7}{g_0^8}.
\end{equation}
\end{corollary}

\begin{proof}
对 $0<c<g_0$，置
\[
 N=256(g_0-c),\qquad K=H_c/N.
\]
命题~\ref{gr:source} 给出 $K\in L^1$ 及
$LK=N^{-1}(\mu_c-\delta_0)$，且 $K\ge256^{-1}\ind_Q$ 几乎处处。
核的零集修改不影响其与 Lebesgue 密度的任意平移、缩放卷积，故此本质下界已经给出逐尺度的立方体平均支配。
由定理~\ref{gr:comparison} 及 \eqref{gr:mass-formula}，
\[
 C_8\le\int K
 =\frac{7J}{8\cdot256}\frac{c^{-1/7}}{g_0-c}.
\]
右侧关于 $c$ 的对数导数为 $(8c-g_0)/(7c(g_0-c))$，唯一极小点为 $c=g_0/8$，代入即得 \eqref{gr:conditional-J}。
在所述对称性下，十六个面的积分相同，每面的 $128$ 个正交象限积分也相同，故 $J=16\cdot128I=2048I$，进而得到 \eqref{gr:conditional-I}。
\end{proof}

\begin{remark}
以上比较使用同一半径作用于八个坐标，且固定 Lebesgue 背景。
任意有限方向组未必具有推论中的反射与面传递对称性，此时应保留总面积分 $J$。
式 \eqref{gr:conditional-I} 只有在同一真实模型的下界 $g_0\le g$ 与上界 $I\le I_0$ 都成立时才能代入。
\end{remark}

\subsubsection{固定面上的短时间与长时间估计}

为了在面上估计 $G$，还需要控制时间积分的两端。
下面的短时间估计覆盖整个平面，并包括奇异子空间的邻域；长时间则分别给出有限面域和全平面的估计。

令
\[
 J_1=\{j:(v_j)_1\ne0\},\qquad
 D_1=\sum_{j\in J_1}|(v_j)_1|,\qquad
 A=\sum_{j=1}^mc_j|(v_j)_1|.
\]
张成性保证 $D_1,A>0$。
对每个 $j\in J_1$ 取包含 $v_j$ 的方向基 $B_j$，并置
\begin{equation}\label{gr:face-B}
 B_* =D_1^2\sum_{j\in J_1}
 \frac{c_j}{\pi^8|\det B_j|\prod_{i\in B_j\setminus\{j\}}c_i}.
\end{equation}

\begin{proposition}\label{gr:face-estimates}
对每个 $t>0$，固定面 $x_1=1$ 上有
\begin{equation}\label{gr:face-pointwise}
 \sup_{y\in\R^7}p_t(1,y)\le B_*t^{-6},\qquad
 \int_{\R^7}p_t(1,y)\,\dd y
 \le\frac{At}{\pi(1+A^2t^2)}.
\end{equation}
记
$T_\varepsilon(x)=\int_0^\varepsilon p_t(x)\,\dd t$，
$R_T(x)=\int_T^\infty p_t(x)\,\dd t$。则
\begin{equation}\label{gr:short-time}
 \|T_\varepsilon(1,\cdot)\|_{L^{8/7}(\R^7)}
 \le\frac89B_*^{1/8}(A/\pi)^{7/8}\varepsilon^{9/8}.
\end{equation}
对任意有限七维测度的 $E\subset\R^7$，
\begin{equation}\label{gr:long-compact}
 0\le R_T(x)\le\frac C{7T^7},\qquad
 \|R_T(1,\cdot)\|_{L^{8/7}(E)}
 \le |E|^{7/8}\frac C{7T^7}.
\end{equation}
全平面另有
\begin{equation}\label{gr:long-plane}
 \|R_T(1,\cdot)\|_{L^{8/7}(\R^7)}
 \le\frac87C^{1/8}(\pi A)^{-7/8}T^{-7/8}.
\end{equation}
\end{proposition}

\begin{proof}
固定 $0<\beta<1$。若所有 $j\in J_1$ 的系数均满足 $|s_j|\le\beta/D_1$，则
$|\sum_js_j(v_j)_1|\le\beta$。
所以落入开区域 $|x_1|>\beta$ 的随机和中，至少一个法向非零的系数满足 $|s_j|>\beta/D_1$。
将该系数限制在这一集合上，其未归一化密度至多为
$c_jtD_1^2/(\pi\beta^2)$；基 $B_j$ 中其余七项的密度至多为 $(\pi c_it)^{-1}$。
基变换、与其余概率律卷积及事件并集给出
\[
 p_t(x)\le\beta^{-2}B_*t^{-6}
 \quad\hbox{几乎处处于 }\{|x_1|>\beta\}.
\]
先由密度连续性得到此开区域的逐点界，再在 $x_1=1$ 处令 $\beta\uparrow1$，得到第一个估计。

第一坐标的边缘律是尺度 $At$ 的一维 Cauchy 律。
边缘密度的识别直接给出的只是对几乎处处 $s$ 成立的等式
\[
 \int_{\R^7}p_t(s,y)\,\dd y=\frac{At}{\pi((At)^2+s^2)}.
\]
固定 $t$，取满足此等式的 $s_k\to1$，利用点态连续性及 Fatou 引理，才得到
\[
 \int_{\R^7}p_t(1,y)\,\dd y
 \le\liminf_{k\to\infty}\int_{\R^7}p_t(s_k,y)\,\dd y
 =\frac{At}{\pi(1+A^2t^2)}
 \le\min\left\{\frac A\pi t,\frac1{\pi At}\right\}.
\]
因此固定切片上使用的是上述方向的不等式。

对非负函数 $h$，$\|h\|_{8/7}\le\|h\|_\infty^{1/8}\|h\|_1^{7/8}$。
由 \eqref{gr:face-pointwise} 的短时间部分，
\[
 \|p_t(1,\cdot)\|_{8/7}
 \le B_*^{1/8}(A/\pi)^{7/8}t^{1/8},
\]
其中时间幂为 $-6/8+7/8=1/8$。
先在有限时间与有限空间区域使用 Minkowski 不等式，再取单调极限，得到
\[
 \|T_\varepsilon(1,\cdot)\|_{8/7}
 \le B_*^{1/8}(A/\pi)^{7/8}\int_0^\varepsilon t^{1/8}\,\dd t,
\]
即 \eqref{gr:short-time}。
$S_3\cap\{x_1=1\}$ 为七维零测集，其无穷点值与此有限范数相容。
若模型另外具有各自由坐标的反射对称性，正象限范数可再乘 $2^{-49/8}$；限制到单位正面仍保留该上界。

积分全空间峰值 \eqref{gr:global-density} 即得 \eqref{gr:long-compact}。
全平面不能直接积分一个统一正常数；这里结合 $\|p_t(1,\cdot)\|_1\le(\pi At)^{-1}$，另得
\[
 \|p_t(1,\cdot)\|_{8/7}
 \le C^{1/8}(\pi A)^{-7/8}t^{-15/8}.
\]
对 $t\in[T,\infty)$ 使用 Minkowski 不等式并积分，便得到 \eqref{gr:long-plane}。
\end{proof}

\section{应用}\subsection{水平集估计与强型推论}
定理\ref{thm:main}的上界对任意输入给出显式分布估计。记
\[
 U=\frac{56627260639}{2000000000}.
\]
则对 $f\in L^1(\R^8)$、$\lambda>0$，
\begin{equation}\label{app:distribution}
 |\{M_8f>\lambda\}|\le \frac{U}{\lambda}\|f\|_1.
\end{equation}
下界构造同时说明，若要求式\eqref{app:distribution}对全部 $L^1$ 输入成立，右侧的统一常数不能小于
\[\frac{2188811179426669}{10^{15}}.\]

\begin{corollary}\label{app:Lp}
对每个 $1<p<\infty$，有
\[
 \|M_8f\|_p\le\frac{p}{p-1}(pU)^{1/p}\|f\|_p,
 \qquad f\in L^p(\R^8).
\]
\end{corollary}
\begin{proof}
固定 $0<\eta<1$。把 $|f|$ 分成
$g=|f|\ind_{\{|f|>\eta\lambda\}}$ 与余项 $h$。此时 $g\in L^1$，$M_8h\le\eta\lambda$，从而
\[
 |\{M_8f>\lambda\}|\le
 \frac{U}{(1-\eta)\lambda}\int_{\{|f|>\eta\lambda\}}|f|.
\]
乘以 $p\lambda^{p-1}$ 后积分，由 Tonelli 得
\[
 \|M_8f\|_p^p\le
 \frac{pU}{(p-1)(1-\eta)\eta^{p-1}}\|f\|_p^p.
\]
取 $\eta=(p-1)/p$ 并开 $p$ 次根即得结论。这是由显式弱型常数得到的直接推论，不涉及强型常数的最优性。
\end{proof}

\subsection{十六方向模型的精确点值与长时间估计}
令
\[
 W=\{(1,\varepsilon_2,\varepsilon_3,\varepsilon_4):
 \varepsilon_j\in\{-1,1\}\},\qquad
 V=(W\times\{0\}^4)\cup(\{0\}^4\times W),
\]
全部方向权取 $1$。对应过程的两个四维分量独立，且使用同一时间。因此
\[
 p_t^{(8)}(x,y)=p_t^{(4)}(x)p_t^{(4)}(y),\qquad
 G_8(x,y)=\int_0^\infty p_t^{(4)}(x)p_t^{(4)}(y)\dd t.
\]
该模型具有全部坐标反射对称性，块内置换与交换两块的群在八个坐标上可迁，故第三节的单面质量公式适用。

取 $A=(1,1,1,0)$、$B=(1,1,0,0)$。四维 Fourier 空间可分成 $1152$ 个单纯锥；逐锥积分并合并全和，可写成
\[
 p_t^{(4)}(A)=\pi^{-4}P_A(t),\qquad
 p_t^{(4)}(B)=\pi^{-4}P_B(t),
\]
其中两者均为准确有理函数。特别
\begin{equation}\label{app:PB}
 P_B(t)=\frac{25344t^4+250t^2+1}
 {7077888t^8+860160t^6+33984t^4+540t^2+3}.
\end{equation}
$P_A$ 的系数及完整锥恒等式见补充材料。对有理函数全和先约分，再对乘积作部分分式分解，利用
\[
 \int_0^\infty(t^2+a^2)^{-k}\dd t
 =\frac{\pi\binom{2k-2}{k-1}}{2\,4^{k-1}a^{2k-1}}
 \quad(a>0,\ k\in\N)
\]
得到下表的准确点值。
\begin{center}
\begin{tabular}{cc}
\toprule
$x$ & $\pi^7G_8(x)$\\
\midrule
$(A,A)$ & $10495703317/2149908480000$\\
$(A,B)$ & $15547486331/1920360960000$\\
$(B,B)$ & $143/10584$\\
\bottomrule
\end{tabular}
\end{center}
这些点的全有理密度在 $t=0$ 有界，乘积在无穷远为 $O(t^{-8})$，故积分恒等式的两端均有限。

有理函数的最高次项给
\[
 \lim_{t\to\infty}t^4P_A(t)=\lim_{t\to\infty}t^4P_B(t)
 =\frac{11}{3072}.
\]
由密度缩放及正时间连续性，此极限就是 $\pi^4p_1^{(4)}(0)$。Fourier 反演的非负乘子进一步给
\[
 0<p_t^{(4)}(x)\le p_t^{(4)}(0)
 =\frac{11}{3072\pi^4}t^{-4}.
\]
从而八维模型满足全空间估计
\begin{equation}\label{app:exact-tail}
 p_t^{(8)}(x)\le\frac{121}{9437184\pi^8}t^{-8},\qquad
 \int_T^\infty p_t^{(8)}(x)\dd t
 \le\frac{121}{66060288\pi^8}T^{-7}.
\end{equation}
它相对于只用八方向基所得的峰常数缩紧 $36864/121$ 倍。在固定面 $x_1=1$ 上，另可取第三节的 $A=8$、$B=2/\pi^8$ 于短时估计。这些显式量可用于截时积分的误差控制。

上述点值给任何允许的立方体下界参数一个必要约束
$g_0\le G_8(A,A)$，并可校验 Green 求值程序。构造新的势核上界还需要整个方体的 $G\ge g_0$ 以及完整面的 $G^{8/7}$ 积分上界；三个点值和式\eqref{app:exact-tail}尚不提供这两项结论。

\subsection{两个后续问题}
几何上界中的角帽积分忽略了不同立方体中心之间的一部分共同选择约束。单独改进每个增量点的中心可行性不能改变最坏角点的局部积分，因此进一步改善应使用跨立方体关系，并在一般 $L^1$ 水平集上重新完成覆盖估计。方向势方法则把问题转化为连续的 $g_0$ 与面积分认证；八维的奇异集应按低秩子空间处理，第三节的整体短时范数可避免逐个切除其交叠邻域。这两个问题分别涉及几何选择和连续核认证，与有限见证族的最优参数问题不同。

\appendix\section{准确计算与补充材料}\subsection{下界的有限证书}
补充材料 \path{paper_claims.json} 保存本稿的准确输入和全分数。
令 $D=10^{35}$、$S=4D$。质量根以式\eqref{pre:root}隔离；事件根 $1$ 取准确值，$1944^{1/8}$ 取相邻的 $D^{-1}$ 有理端点。所有 $32768$ 个标签都重建为固定坐标次序的物理盒，再分别求内、外并集。下表记录足够复核本下界的有限规模。
\begin{center}
\begin{tabular}{lr}
\toprule
项目&数量\\\midrule
带标签源块&$32768$\\
正内盒、正外盒&各 $882$\\
去除完全重复后的物理内盒、外盒&各 $644$\\
准确并集的总缓存状态&$210446$\\
共同半径的全部源坐标包含检查&$34054$\\
\bottomrule
\end{tabular}
\end{center}
完整有理区间位于
\[
\begin{gathered}
 L_0\le 2^8\left|\bigcup B_\ell^-\right|
 \le 2^8\left|\bigcup B_\ell^+\right|\le U_0,\\
 L_0=\frac{21888111794266698273227523980123024}{10^{34}},\\
 U_0=\frac{21888111794266698273227523980123027}{10^{34}}.
\end{gathered}
\]
这里两端是由原始全分数直接向外舍入所得的简化有理界；验证程序仍比较完整内外分数。
取正文中的 $\varepsilon,\theta,K$ 后，程序验证
\[
 \theta\left(L_0-\frac{441}{9536743164062500000}\right)
 \left(\frac{2K+1}{2K+5}\right)^8
 >\frac{2188811179426669}{10^{15}}.
\]

\subsection{上界标量的有理判据}
取 $t_0=3837029/524288$，把 $B_8(t_0)$ 按二项式展开写成
$\alpha+\beta\log t_0$，其中 $\alpha,\beta\in\mathbb Q$。
令 $z=(t_0-1)/(t_0+1)$。使用 $N=400$ 时，
\[
 0\le\log t_0-2\sum_{j=0}^{N-1}\frac{z^{2j+1}}{2j+1}
 \le\frac{2z^{2N+1}}{(2N+1)(1-z^2)}.
\]
按 $\beta$ 的符号传播上下端，得到 $B_8(t_0)\le B_+$，其中可取
\[
 B_+=\frac{133365553664751128154150225918012012532222}{10^{40}}.
\]
记 $s=(1+1/t_0)^4$、$U=56627260639/2000000000$。准确计算满足
\[
 d:=U-s^2-B_+>0,\qquad
 d^2-4s^2B_+>\frac{19}{10^{10}}>0.
\]
因为 $B_+>0$，两次平方保序，由此得到
$F_8(t_0)\le(s+\sqrt{B_+})^2<U$。

\subsection{四维密度的完整锥分解}
为说明应用中点值计算的有限化，取频率排序室
$\xi=(a,b,c,d)$，$a\ge b\ge c\ge d\ge0$。定义射线
\[
\begin{aligned}
 r_1&=(1,0,0,0),&r_2&=(1,1,0,0),&r_3&=(1,1,1,0),\\
 r_4&=(1,1,1,1),&r_5&=(2,1,1,0),&r_6&=(3,1,1,1).
\end{aligned}
\]
排序室按 $a\ge b+c+d$、$a\le b+c-d$ 和中间区域分为
\[
 C_1=\operatorname{cone}(r_1,r_2,r_5,r_6),\quad
 C_2=\operatorname{cone}(r_5,r_2,r_3,r_4),\quad
 C_3=\operatorname{cone}(r_5,r_2,r_4,r_6).
\]
三组锥坐标分别为
\[
\begin{gathered}
 (a-b-c-d,\ b-c,\ c-d,\ d),\\
 (a-b,\ b-c,\ -a+b+c-d,\ d),\\
 \left(c-d,\ b-c,\ \frac{-a+b+c+d}{2},\
                  \frac{a-b-c+d}{2}\right).
\end{gathered}
\]
在相应区域它们非负，且准确重构频率向量，证明三锥覆盖及内部不交。
绝对行列式依次为 $1,1,2$。
对 $\psi_4(\xi)=\sum_{v\in W}|v\cdot\xi|$，六条射线的值依次为
$8,8,12,12,16,24$；每条 $v$ 在一个锥的全部射线上同弱号，故
$\psi_4$ 在该锥上线性。
全部 $2^4\,4!=384$ 个带符号置换给出 $1152$ 个锥，交叠只在零测边界。
因此 Fourier 反演逐锥积分精确给
\[
 p_t^{(4)}(x)=\frac1{16\pi^4}
 \sum_C\frac{|\det(r_{C,1},\ldots,r_{C,4})|}
 {\prod_{j=1}^4\bigl(t\psi_4(r_{C,j})-i x\cdot r_{C,j}\bigr)}.
\]
$t>0$ 时分母均不为零。对 $x=A,B$，将全部项通分后再约分，得到补充数据中的 $P_A,P_B$。
可选程序 \path{supplement/verify_green.py} 重新生成全部锥，对实、虚分子作准确全和，验证约分后的有理函数及三个部分分式积分。零投影和重复投影均保留在全和中处理。

\subsection{复现与材料范围}
\path{supplement/verify_paper.py} 仅依赖 Python 标准库，重建并复算本稿唯一的下界构造，核对有限光滑输入的比值，并以有理数验证上述对数尾和平方判据。它不执行搜索，不计算其他首轴参数，不认证某个构造族的全局最优值。
\begin{verbatim}
python3 -B supplement/verify_paper.py
\end{verbatim}
较快的 \texttt{--scalar-only} 模式只验证保存分数与上界标量，不代替完整并集复算。
\path{green_anchors.json} 保存两个四维有理密度和三个准确积分的部分分式。补充说明另外标明原始程序、来源散列与所需软件；主定理的标准库复核不依赖浮点优化器、MATLAB 或区间库。

正文中的主定理、条件势核推论与点值计算具有不同作用：主定理已给出 $C_8$ 的双侧界；条件势核推论需另有全域 $g_0,I_0$ 证书；点值仅提供校验与必要限制。研究过程记录、版本散列和执行回执不作为额外数学假设。

\begin{thebibliography}{99}
\bibitem{Aldaz2000} J. M. Aldaz. A remark on the centered $n$-dimensional Hardy--Littlewood maximal function. \emph{Czechoslovak Mathematical Journal}, 50 (2000), 103--112. \url{https://dml.cz/handle/10338.dmlcz/127554}.
\bibitem{SteinStromberg1983} E. M. Stein and J.-O. Str\"omberg. Behavior of maximal functions in $\R^n$ for large $n$. \emph{Arkiv f\"or Matematik}, 21 (1983), 259--269. \url{https://doi.org/10.1007/BF02384314}.
\bibitem{Aldaz2011} J. M. Aldaz. The weak type $(1,1)$ bounds for the maximal function associated to cubes grow to infinity with the dimension. \emph{Annals of Mathematics}, 173 (2011), 1013--1023. \url{https://doi.org/10.4007/annals.2011.173.2.10}.
\bibitem{IakovlevStromberg2013} A. S. Iakovlev and J.-O. Str\"omberg. Lower bounds for the weak type $(1,1)$ estimate for the maximal function associated to cubes in high dimensions. \emph{Mathematical Research Letters}, 20 (2013), 907--918. \url{https://doi.org/10.4310/MRL.2013.v20.n5.a7}.
\bibitem{AldazKissing2020} J. M. Aldaz. Kissing numbers and the centered maximal operator. arXiv:1811.10478v2, 2020. \url{https://arxiv.org/abs/1811.10478v2}.
\bibitem{Lin2026} Y. Lin. Two partial balayage principles for weak-type estimates. 网页手稿，访问日期2026年10月27日. \url{https://coolrmal.github.io/articles/two-partial-balayage-principles/}.
\bibitem{Folland1999} G. B. Folland. \emph{Real Analysis: Modern Techniques and Their Applications}. Second edition. John Wiley \& Sons, 1999.
\bibitem{Evans2010} L. C. Evans. \emph{Partial Differential Equations}. Second edition. Graduate Studies in Mathematics 19, American Mathematical Society, 2010.
\end{thebibliography}

\end{document}
